\section{Introduction}
The \textit{Online Packet Scheduling with Deadlines} (OPSD, for short) problem, also called \textit{bounded delay buffer problem} or \textit{traffic shaping}, is a theoretical framework that has been introduced to model the behavior of buffers that store data packets in \textit{Quality-of-Service} (QoS) networks \citep{karakus2017quality}. In QoS networks, it is crucial to enforce service standards, especially on the performance as seen by the end-users. QoS plays a key role in the modern digital infrastructure, notable examples include sensor networks \citep{iyer2003qos}, cloud computing \citep{szigeti2013end} and Intelligent Transportation Systems (ITS) \citep{belamri2021survey}.

In its simpler version, the OPSD framework assumes that packets with unit processing cost arrive over time and are stored in an unlimited buffer. 
%A packet $p$ is characterized by an arrival time $r_p$, a \textit{deadline} $d_p$, and a \textit{weight}. At every time unit $t$ (\emph{e.g.}, a clock cycle), a \textit{scheduler} is allowed to process one packet $p_t$ from the buffer, \emph{i.e.}, the arrival time must satisfy $r_{p_t} \le t$. 
In OPSD, the presence of \textit{deadlines} models the practical scenario of low-latency systems. In other words, it is not possible to keep a packet idling in the buffer for too many clock cycles. Every packet must then either be processed before its requested deadline or discarded from the buffer. The model is completed by including scores representing the importance of each packet, called \textit{weights}. Of course, any algorithm \alg has the goal of maximizing its total score, or \textit{cumulative gain}. 
%To simplify the notation, we directly indicate the weight of a packet with its name $p$. 
Naturally, the presence of the weights makes the problem non-trivial and models the QoS property of the system. The arrival sequence defines an \textit{instance} of the OPSD problem and can be chosen in an arbitrary way by an \textit{adversary}. As different instances may present different opportunities to \alg, a standard way to evaluate its performance is to compare it to the best possible decision-maker for that instance---called \opt---through the \textit{competitive ratio} $C$, which is the ratio between the gain of \opt and the gain of \alg.

% The goal of a scheduler is to maximize its total gain, defined as the sum of the weights of the scheduled packets. For a general scheduling algorithm \alg, we call $\Galg = \sum_{t} p_t$ its cumulative gain, and we compare it with \opt, the best possible scheduler (aware of the full arrival sequence beforehand) with total gain $\Gopt \ge \Galg$ for every \alg. A natural performance metric is the \textit{competitive ratio}, defined as the smallest value $C \ge 1$ s.t. $\Gopt \le C\cdot \Galg$ for every instance. 

% The literature on OPSD can be segmented based on two major distinctions: \textit{deterministic} and \textit{randomized} algorithms and the \textit{slackness degree}.

% We say that a system allows for randomization if the algorithm is allowed to undertake randomized decisions without disclosing the state of its \textit{random bits}, \emph{i.e.}, the adversary can fully observe the algorithm and the probabilities with which it takes a decision beforehand, but not the actual realizations. If the adversary is allowed to completely observe the future decisions of the algorithm, then we say that it is deterministic.

% An OPSD instance can be characterized by the maximum amount of \textit{slackness} of its packets. The slackness is the difference between the deadline and the arrival time of a packet. In $s$-bounded OPSD instances, with $s \ge 1$, the maximum slackness is $s-1$. In the degenerate case of $s=1$, all the packets expire in the same time slot as they appear. 
% When $s=\infty$, there is no known upper bound on the slackness. The OPSD literature treats settings with different degrees of slackness separately. A particularly relevant scenario is the $2$-bounded OPSD problem, where the amount of slackness is minimal, but the problem remains non-trivial.

Over the last few years, there has been a surge in interest in \textit{learning-augmented} online algorithms. The idea is to include \textit{online learning techniques} to deal with reduced information or to boost the performance of online algorithms. Examples of online optimization problems with predictions include OPSD \citep{liang2023learning}, facility location \citep{agrawal2022learning}, and paging \citep{bansal2022learning}. Our work initiates the study of a variant of the OPSD problem, which we call \textit{Stochastic Online K-Packet Scheduling with Deadlines} ($K$-OPSD, for short), where we include online learning (and, in particular, \textit{bandit learning}) into the standard OPSD problem. In the $K$-OPSD framework, we assume that every packet belongs to one of $K$ different types, which is known by the scheduler. However, the weight of a packet is not revealed until it is processed. We make a stochastic assumption, assuming that every packet type is associated with a probability distribution, and when a packet is processed, the weight is sampled from the corresponding distribution. Recently, this type of extension appeared in \cite{merlis2023preemption} in the context of preemptive scheduling and in \cite{levy2024online} for online paging.

% The unknown and stochastic weights assumption models real-world applications where the weight of a packet is hidden in the metadata due to encryption or privacy constraints, and a packet weight can only be inferred after execution. Also, this is a natural model for QoS networks when the packets to be executed are user-level requests, and the goal is to execute requests that will result in more downstream usage. Again, the weight can only be inferred after execution, depending on the user's type. 
The goal of a scheduler in the $K$-OPSD problem can be naturally extended from the classic OPSD setting: a scheduling algorithm \alg wants to maximize its \textit{expected} cumulative gain.
%, defined as $\mathbb{E}[\Galg] = \sum_{t}p_t$. In the same way, we define an optimal scheduler \opt that proposes a schedule that is \textit{optimal in expectation}, that implies $\mathbb{E}[\Gopt] \ge \mathbb{E}[\Galg]$ for every \alg. Intuitively, \opt observes the whole arrival sequence beforehand, does \textit{not} observe the actual weight realizations, but rather their expectations. 
In $K$-OPSD, a scheduling algorithm has less information available than in OPSD. As learning becomes a necessary component for a scheduler, a natural way to evaluate its performance is $\alpha$-regret. The notion of $\alpha$-regret combines the competitive ratio with an additive regret w.r.t. an optimal decision-maker that knows the average weights for every packet type and does not need to learn them. 
%For this definition, the number of available samples becomes relevant, and we call $T$ the \textit{time horizon} of the problem. $T$ corresponds to the length of the schedule that \alg is requested to compute. We define the $\alpha$-regret as $R_{\alpha,T}^{\text{\alg}} \coloneqq \mathbb{E}[\Gopt] - \alpha\mathbb{E}[\Galg]$, for some $\alpha \ge 1$. The smaller $\alpha$ and $R_{\alpha,T}^{\text{\alg}}$ are, the better the metric. 

% The notion of regret has been extensively studied in the Multi-Armed Bandit literature (MAB, \cite{lattimore2020bandit}). In the stochastic MAB problem, a learner \alg is faced with $K$ possible actions. At every round $t \le T$, the learner has to choose one of the actions $I_t$ and receive a reward sampled from a fixed and unknown distribution $\mathcal{D}_{I_t}$, with expectation $\mu_{I_t}$. The goal is to maximize the cumulative expected reward $\mathbb{E}[\Galg] = \sum_{t=1}^T \mu_{I_t}$. \textit{The expected regret} of \alg is computed as $R_T^{\text{\alg}} \coloneqq T\max_{i\in[K]}\mu_{i}-\mathbb{E}[\Galg]$, as the performance gap with the best possible decision maker in expectation. In the MAB literature, the goal is usually to provide an algorithm for which the expected cumulative regret can be upper bounded sub-linearly in $T$, \emph{i.e.}, $R_T^{\text{\alg}} = o(T)$. Algorithms with sub-linear regret guarantees are called \textit{no-regret} algorithms. We call $\alpha$-no regret algorithms the ones having sub-linear $\alpha$-regret, \emph{i.e.}, $R_{\alpha,T}^{\text{\alg}} = o(T)$.

The stochastic MAB problem \cite{lattimore2020bandit} shares many similarities with the $K$-OPSD problem. 
We show that the $K$-OPSD problem includes (and generalizes) the Sleeping Bandit problem \cite{kleinberg2010regret}, which is a generalization of the classic MAB problem where an adversary can choose the subset of available actions at every round.
% In fact, the $K$-OPSD setting is a generalization of the stochastic MAB problem, where every packet type can be mapped to an action, and where the set of available actions may change depending on past decisions.
The OPSD setting is included in the $K$-OPSD setting, and even in the simpler scenarios, it is possible to prove a lower bound on the competitive ratio of any algorithm. For example, even in $2$\textit{-bounded instances} where the deadline of every packet is upper bounded by the arrival time plus one, a lower bound on the competitive ratio is $\Phi = \frac{1+\sqrt{5}}{2} \approx 1.618$, the so-called \textit{golden ratio}. A competitive ratio greater than one implies the impossibility of obtaining a sub-linear \emph{standard} regret in the $K$-OPSD problem. In this paper, we address the following question:
\begin{center}
    \textit{Is it possible to devise $\alpha$-no-regret algorithms in some $K$-OPSD problems?} 
\end{center}
We answer this question positively. It turns out that in many relevant scenarios this is possible, and with provably tight regret rates and state-of-the-art competitive ratios. Interestingly, the $K$ types assumption can even be leveraged to improve the existing competitive ratios.
\paragraph{Motivating Example}
As previously discussed, the OPSD problem is well-motivated by several applications. We now state a real-world scenario in which $K$-OPSD is a better fit, which motivates the setting. We are in a digital advertising scenario, in which there exist $K$ ad publishers that, in a really fast-paced manner, ask to participate in an auction to display an ad on their space. As an advertiser, you are interested in participating in the auctions with the highest possible probability of receiving a click, which is, however, unknown. However, the requests are often so large in number that not all of them can be processed. If you don't participate in the auction, you will never know if you would have received a click or not. Moreover, auctions have a limited lifespan and usually finish after a few seconds. This problem combines reward estimation and scheduling with deadlines, since a quick decision based on past observation must be made. Moreover, as the lifespan of auctions is usually \textit{really} short, it is even a stronger motivation for scenarios with short deadlines, such as the $2$-bounded $K$-OPSD problem.

\subsection{Existing Results and Summary of Contributions}

\begin{table*}[t!]
\centering
\setlength{\tabcolsep}{5pt}
\renewcommand{\arraystretch}{1.2}
\begin{tabular}{llll}
\toprule
\textbf{Setting} & \makecell{\textbf{Algorithm} \\ \textbf{Type}} & \makecell{\textbf{OPSD} \\ \textbf{Competitive Ratio} }&  \makecell{\textbf{$K$-OPSD}  \\ \textbf{$\alpha$-regret} }\\
\midrule
\multirow{2}{*}{\makecell[l]{$1$-bounded \\ (MAB)}}
  &Det. (upper)          & $1$  & $R_{1,T} \le \widetilde{\mathcal{O}}\left(\sqrt{KT}\right)$ 
  ~\cite{kleinberg2010regret}           \\
  \addlinespace[-2pt]
  & Det. (lower)           & $1$  & $R_{1,T} \ge \Omega\left(\sqrt{KT}\right)$ ~\cite{kleinberg2010regret}\\
\midrule
\multirow{4}{*}{$2$-bounded}
%   & \multirow{2}{*}{Det. (upper)}           & $\Phi$~~~~\cite{kesselman2001buffer}   & 
%   \multirowcell{2}[-0.5ex][l]{\cellcolor{gray!20}$R_{\theta_K, T} \le \widetilde{\mathcal{O}}\bigl(\sqrt{KT}\bigr),~~\theta_K < \Phi$}
% \\
%   & & \cellcolor{gray!20}$\theta_K \stackrel{K \rightarrow \infty}{\rightarrow}\Phi$ & \\
  & Det. (upper) & \cellcolor{gray!20}$\theta_K \stackrel{K \rightarrow \infty}{\rightarrow}\Phi, ~~  (\Phi$~~\cite{kesselman2001buffer})  & \cellcolor{gray!20}$R_{\theta_K, T} \le \widetilde{\mathcal{O}}\bigl(\sqrt{KT}\bigr),~~~\theta_K < \Phi$ \\
  & Det. (lower)          & $\theta_K \stackrel{K \rightarrow \infty}{\rightarrow}\Phi$~~~~\cite{hajek2001competitiveness} & \cellcolor{gray!20} {$R_{\theta_K, T} \ge \Omega\left(\sqrt{T}\right),~~~~~\theta_{K-1} < \Phi$} \\
  & Rand. (upper)   & $\frac{5}{4} $~~~~\cite{chin2006online}    & \cellcolor{gray!20} $R_{\frac{5}{4}, T} ~\le \widetilde{\mathcal{O}}\left(\sqrt{KT}\right)  $             \\
  & Rand. (lower)   &$ \frac{5}{4} $~~~~\cite{chin2003online}     & N.A.                \\
\midrule
\multirow{2}{*}{$3$-bounded}
  & Det. (upper)           & $\Phi$~~~~\cite{kesselman2001buffer}   & \cellcolor{gray!20} $R_{\Phi, T} ~\le \widetilde{\mathcal{O}}\left(\sqrt{KT}\right)$              \\
  & Det. (lower)           & $\Phi$ & \cellcolor{gray!20} $R_{\theta_K, T} \ge \Omega\left(\sqrt{T}\right),~~~\theta_{K-1} < \Phi$ \\
\midrule
\multirow{4}{*}{\makecell[l]{Unbounded \\ Slackness}}
& Det. (upper)           & $\Phi$~~~~\cite{vesely2022competitive} &   N.A.        \\
& Det. (lower)           & $\Phi$  & \cellcolor{gray!20} $R_{\theta_K, T} ~\ge \Omega\left(\sqrt{T}\right),~~\theta_{K-1} < \Phi$        \\
& Rand. (upper)   & $\frac{e}{e-1}$~~\cite{chin2006online}  &  \cellcolor{gray!20} $R_{\frac{e}{e-1}, T} \le \widetilde{\mathcal{O}}\left(\sqrt{KT}\right)$  \\
& Rand. (lower)   & $\frac{5}{4}$ &   N.A. \\
\bottomrule
\end{tabular}
\caption{Comparison of existing results with ours. Results without citations are direct consequences of other entries. Results highlighted in gray are our contributions. The $\widetilde{\mathcal{O}}$ notation hides universal constants and logarithmic factors. The $\Omega$ notation hides universal constants.}
\label{tab:literature}
\end{table*}

In Table \ref{tab:literature}, we summarize the existing results in the classic OPSD problem and compare them to our results. The $K$-OPSD setting is novel and unexplored in the literature, thus our bounds on the $\alpha$-regret can only be compared in light of the existing competitive ratios and upper bounds on the regret. However, in the $2$-bounded setting, we improve over existing results when $K<\infty$. As customary in this literature, we distinguish between \textit{deterministic} and \textit{randomized} algorithms. In what follows, we report an overview of our results. Through the paper, we use the $\mathcal{O}$ notation to hide universal constants, and the $\widetilde{\mathcal{O}}$ notation to hide universal constants and logarithms.

In the first part of this paper, we propose deterministic algorithms for the OPSD and the $K$-OPSD setting.
\begin{itemize}
    \item In the $2$-bounded and $3$-bounded $K$-OPSD setting, we propose a deterministic algorithm that suffers a $\Phi$-regret upper bounded as $\widetilde{\mathcal{O}}\left(\sqrt{KT}\right)$ (Theorem \ref{prop:reduction}). This upper bound is also provided in an instance-dependent version. In classic Sleeping MABs the best attainable bound on the the regret is in the order of $\widetilde{\mathcal{O}}\left(\sqrt{KT}\right)$ as well, and the competitive ratio matches the best existing one. This result highlights the connection between the $K$-OPSD setting and the Sleeping MAB setting, and the analysis provides insights into how the learning procedure can be decomposed in a meaningful way from the traditional competitive ratio analysis.
    \item In the $2$-bounded OPSD setting, when the number of distinct weights is $K<\infty$, we provide a novel deterministic algorithm called \algtheta that achieves a competitive ratio of $\theta_K$ (\ref{prop:upperboundtheta}) smaller than the best known competitive ratio of $\Phi$ \cite{hajek2001competitiveness}. Note that $\theta_K \in [\sqrt{2}, \Phi)$ for $K<\infty$ and $\theta_\infty = \Phi$. This upper bound on the competitive ratio is tight according to the lower bound proposed in \cite{hajek2001competitiveness}.
    \item In the $2$-bounded $K$-OPSD setting, we adapt \algtheta to deal with learning, resulting in \algthetalearn. This deterministic algorithm suffers a $\theta_K$-regret bounded as $\widetilde{\mathcal{O}}\left(\sqrt{KT}\right)$ (Theorem \ref{thr:upperbound_thetalearn}), and we prove that this is nearly tight with a lower bound (Theorem \ref{thr:lowerbound}).
\end{itemize}

We dedicate the second part of this paper to randomized algorithms.
\begin{itemize}
    \item  In the $2$-bounded $K$-OPSD setting, we propose \algrtwo, a randomized algorithm that suffers a $\frac{5}{4}$-regret upper bounded as $\widetilde{\mathcal{O}}\left(\sqrt{KT}\right)$ (Theorem \ref{thr:upperboundrtwo}).
    \item In the $K$-OPSD setting with unbounded slackness, we propose \algrs, an algorithm that suffers a $\frac{e}{e-1}$-regret upper bounded as $\widetilde{\mathcal{O}}\left(\sqrt{KT}\right)$ (Theorem \ref{thr:upperboundrs}).
\end{itemize}
In both cases, our bounds match the best existing competitive ratio as well as the tightest regret attainable in Sleeping MABs.
